\section{Introduction}
\IEEEPARstart{M}{achine} unlearning (MU) has emerged as a critical research area~\cite{hou2025fixguard,ye2023stealthy}, addressing the imperative to prevent unauthorized data exploitation, eliminate harmful model behaviors, and safeguard user privacy~\cite{zhu2024privauditor}. The proliferation of AI misuse ranging from phishing generation and model inversion attacks to perpetuating misinformation has prompted regulatory frameworks~\cite{li2024model,prabhakar2025model,kahla2022label}. Frameworks like the General Data Protection Regulation (GDPR) and California Consumer Privacy Act (CCPA)~\cite{goldman2020introduction,data2018general} mandate that model service providers erase user data from learned representations upon request, establishing legal and ethical obligations for responsible AI deployment. Users possess the right to request data removal at any time, so per-request efficiency becomes a first-class concern.


Existing machine unlearning methods predominantly target image classification~\cite{wuerkaixi2025accurate} and generative tasks~\cite{wei2025emma}, with limited exploration of privacy-critical biometric systems. While face recognition~\cite{de2024unlearning,zakharov2026face,shivam2025cure} has received modest attention, fingerprint and iris recognition remains substantially less studied. This gap persists despite their sensitive nature and widespread deployment in authentication systems. Contemporary approaches rely primarily on KL divergence maximization~\cite{kurmanji2023towards} or gradient ascent~\cite{zhao2024continual}. These methods implement many-to-one mappings that force forgotten class embeddings toward a single distant point in feature space. Alternative strategies include saliency-based selection~\cite{fan2023salun}, geometric boundary manipulation~\cite{chen2023boundary}, perturbation-based forgetting~\cite{sun2024forget}, and meta-learning optimization~\cite{huang2024learning,shi2025mcu}. A key observation across all these methods is that every forget sample is driven toward a single collapse point in the decision space. Even when a sample already lies far from the retain manifold, it must still traverse the full distance to this distant target. This increases convergence time and GPU cost, as observed in Table~\ref{tab:all}. This many-to-one collapse behavior is visualized and ablated in detail in Section~\ref{sec:tsne}. Methods such as GS-LoRA~\cite{zhao2024continual}, SCRUB~\cite{kurmanji2023towards}, SalUn~\cite{fan2023salun}, Boundary Unlearning~\cite{chen2023boundary}, Forget Vectors~\cite{sun2024forget}, MCU~\cite{shi2025mcu}, LTU~\cite{huang2024learning}, DELETE~\cite{zhou2025decoupled}, LoTUS~\cite{spartalis2025lotus}, and MUNBa~\cite{wu2025munba} each trade off convergence speed, memory, or reversibility, and few explicitly target biometric edge deployment or sequential forget cycles.

Biometric recognition differs from general classification in ways that directly affect unlearning feasibility. Biometric systems deploy on resource-constrained edge devices. Removal requests arrive sequentially and frequently over the deployment lifecycle. This transforms unlearning into a recurring operational cost rather than a one-time procedure. Under this workload, existing many-to-one methods become prohibitive as per request compute scales poorly with request frequency. Compounding this, existing unlearning methods have been designed and benchmarked almost exclusively on ResNet~\cite{he2016deep} or ViT backbones~\cite{dosovitskiy2020image} for CIFAR~\cite{krizhevsky2009learning} and ImageNet~\cite{russakovsky2015imagenet}. Their transferability to biometric backbones remains untested, despite architectures such as Face Transformer~\cite{zhong2021face}, LAM~\cite{luo2025novel}, JIPNet~\cite{guan2024joint}, SSLDMNet-Iris~\cite{kadhim2025deep}, IrisFormer~\cite{sun2024irisformer}, PFSimNet~\cite{yu2024partial}, and MEM-FR~\cite{wu2025privacy} representing the state of the art among deep embedding-based biometric backbones. What biometric unlearning requires instead is a cycle-robust method validated directly on biometric backbones.

To address these challenges, we propose Camouflage Unlearning (CU), a novel framework that achieves 3--4$\times$ faster convergence while maintaining robust performance. Unlike existing many-to-one approaches, CU pushes forget-class embeddings toward the nearest retained class boundaries, effectively reducing migration distances. This proximity target governs only the training-time optimization objective, not final verification-level similarity. A naive reading might suggest CU induces impersonation of the nearest retain identity, but Section~\ref{sec:attack} empirically verifies that forget probes do not match their assigned retain centroid more often than chance, ruling out this concern. To perform surgical modifications and minimize collateral damage to retain classes, CU employs localized regularization concentrated in specific blocks of the neural network. Additional details are provided in Section~\ref{sec:4}.

Our main contributions are summarized as follows:
\begin{enumerate}
    \item Among the earliest studies of machine unlearning for privacy-critical deep embedding-based biometric systems, covering fingerprint, iris, and face recognition.
    \item Camouflage Unlearning, achieving $3$--$4\times$ faster convergence than SoTA methods with competitive retention accuracy and privacy guarantees.
    \item Empirical evaluation across biometric and classification domains showing CU concentrates updates in specific network blocks and remains effective under data-scarce replay.
\end{enumerate}
